\emph{Setting.} $S(V)$ is spanned by the monomials $e_{j_1}\cdots e_{j_T}$ in a basis of $V$, ordered, with odd basis elements anticommuting and squaring to zero and even ones commuting; on the even modes it is the palatial polynomial representation (Fact \claimref{field-theory.imported-statistics-no-go-palatial-algebra}) with the positive form $\langle z^n, z^m\rangle = n!\,\delta_{nm}$ stated in the setting of \claimref{synthesis.self-adjoint-brst-charge-needs-a-krein-form-and-the-tower-supplies-the-ghosts}, on the odd modes the exterior algebra of a tower floor (Fact \claimref{field-theory.imported-tower-limit}) with its orthonormal monomials. For a one-particle operator $X$ of parity $|X|$, $d\Gamma(X)$ is the graded derivation of $S(V)$ equal to $X$ on $V$: on a monomial it replaces one factor $e_{j}$ by $Xe_{j}$, with the sign $(-1)^{|X|m}$ for the $m$ odd factors passed. $\Gamma(J)$ is the algebra automorphism equal to $J$ on $V$, and the Fock Krein form is $\langle x, \Gamma(J)y\rangle$, the graded-symmetric extension of $\langle\cdot,\cdot\rangle_K$.

\emph{(i).} The graded commutator $[D_1, D_2] = D_1D_2 - (-1)^{|D_1||D_2|}D_2D_1$ of two graded derivations is a graded derivation, and a graded derivation of $S(V)$ is fixed by its restriction to the generators $V$. On $V$, $[d\Gamma(X), d\Gamma(Y)]$ and $d\Gamma([X, Y])$ agree, so $d\Gamma$ is a homomorphism of Lie superalgebras. Hence $2Q_{B}^2 = [Q_{B}, Q_{B}] = d\Gamma([Q_1, Q_1]) = 2d\Gamma(Q_1^2) = 0$, and $Q_{B}R + RQ_{B} = d\Gamma(Q_1R_1 + R_1Q_1) = d\Gamma(P_{q}) = N_{q}$. Since $V = V_{\mathrm{ph}} \oplus V_{q}$, $S(V) = S(V_{\mathrm{ph}}) \otimes S(V_{q})$ and $N_{q}$ is $n$ on $S(V_{\mathrm{ph}}) \otimes S^{n}(V_{q})$. $N_{q}$ commutes with $Q_{B}$, because $[Q_{B}, Q_{B}R + RQ_{B}] = [Q_{B}^2, R] = 0$, so $Q_{B}$ preserves each eigenspace. On the eigenspace $n > 0$, $Q_{B}x = 0$ gives $x = \tfrac1n(Q_{B}R + RQ_{B})x = Q_{B}(\tfrac1nRx)$: the cohomology vanishes there. On $n = 0$, $Q_{B} = d\Gamma(Q_1)$ acts on $S(V_{\mathrm{ph}})$ through $Q_1V_{\mathrm{ph}} = 0$, so it is zero, and nothing in $S(V_{\mathrm{ph}})$ is exact because $Q_{B}$ preserves $N_{q}$. So $\ker Q_{B} = S(V_{\mathrm{ph}}) \oplus (\ker Q_{B} \cap \{N_{q} > 0\}) = S(V_{\mathrm{ph}}) \oplus \operatorname{im}Q_{B}$ and the cohomology is $S(V_{\mathrm{ph}})$; $V_{\mathrm{ph}}$ carries no ghosts, so it sits at ghost number zero. Krein self-adjointness: $\Gamma(J)$ and the positive form are multiplicative over the monomials, and the Fock adjoint of a one-particle lift is the lift of the one-particle adjoint, so $Q_{B}$ is self-adjoint because $Q_1$ is. Then $\langle x, Q_{B}y\rangle_K = \langle Q_{B}x, y\rangle_K = 0$ for $x \in \ker Q_{B}$: $\operatorname{im}Q_{B}$ is orthogonal to $\ker Q_{B}$ and lies in it, hence is null. The Krein orthogonality of $V_{\mathrm{ph}}$ and $V_{q}$ makes $S(V_{\mathrm{ph}})$ orthogonal to every monomial containing a quartet factor. On $S(V_{\mathrm{ph}})$, with a Krein-orthonormal basis of the positive space $V_{\mathrm{ph}}$, the monomials are orthogonal with norms $\prod n_j! > 0$. So the form on $\ker Q_{B}$ is positive semi-definite with radical $\operatorname{im}Q_{B}$, and positive definite on the cohomology.

\emph{(ii).} With the form $-u\cdot w$ on the polarisation modes (K1) and $\langle c, \bar c\rangle_K = 1$ (K2), self-adjointness of $Q_1$ is $\langle Q_1u, \bar c\rangle_K = -(k\cdot u) = \langle u, k\rangle_K = \langle u, Q_1\bar c\rangle_K$; nilpotency is $Q_1^2\bar c = Q_1k = -Q(k)c = 0$. This is the one-particle form of the charge $L^{\dagger}b + d^{*}L$ of \claimref{synthesis.self-adjoint-brst-charge-needs-a-krein-form-and-the-tower-supplies-the-ghosts}(ii), $c$ and $\bar c$ the one-particle states of its two ghost modes. With $\bar k$ null, $k\cdot\bar k = 1$, and $t$ in the transverse plane $\{k, \bar k\}^{\perp}$: $Q_1\bar k = -c$, $Q_1k = Q_1t = Q_1c = 0$, $Q_1\bar c = k$. For $R_1$ as stated, $(Q_1R_1 + R_1Q_1)$ sends $k \mapsto Q_1\bar c = k$, $\bar k \mapsto R_1(-c) = \bar k$, $c \mapsto Q_1(-\bar k) = c$, $\bar c \mapsto R_1k = \bar c$ and $t \mapsto 0$: it is $P_{q}$. The transverse plane is spacelike, so $-t\cdot t > 0$, and it is Krein-orthogonal to $k$, $\bar k$ (by definition) and to the ghosts. (i) applies; for several colours and momenta, take direct sums of one-particle spaces, charges and homotopies.

\emph{(iii).} Krein self-adjointness with $\langle\phi, \phi\rangle_K = 1$ gives $\langle Q_1u, \bar c\rangle_K = -(p\cdot u) = \langle u, Q_1\bar c\rangle_K$ and $\langle Q_1\phi, \bar c\rangle_K = \beta = \langle\phi, Q_1\bar c\rangle_K$, so $Q_1\bar c = p + \beta\phi$, and $Q_1^2\bar c = (-Q(p) + \beta^2)c$: nilpotency holds exactly when $\beta^2 = Q(p)$. Then $Q_1(p + \beta\phi) = 0$ and $Q_1\phi = \beta c$. Take $R_1c = \beta^{-1}\phi$, $R_1p = \bar c$, and $R_1$ zero on $\phi$, $\bar c$ and $p^{\perp}$. Then $Q_1R_1 + R_1Q_1$ sends $c \mapsto Q_1\beta^{-1}\phi = c$, $\phi \mapsto R_1(\beta c) = \phi$, $p + \beta\phi \mapsto Q_1\bar c = p + \beta\phi$, $\bar c \mapsto R_1(p + \beta\phi) = \bar c$ and $p^{\perp}$ to $0$: it is $P_{q}$ (the check solves the linear system for an odd $R_1$ independently). The complement $p^{\perp}$ is spacelike, Krein-positive and Krein-orthogonal to $p$, $\phi$ and the ghosts, and $Q_1p^{\perp} = 0$; (i) applies. In the field theory the Goldstone shift $s\phi = mc$ and $sA = \nabla c$ give this charge in the Feynman-'t Hooft gauge, where the unphysical modes share the vector's shell $Q(p) = m^2$; in other covariant gauges the quartet lies on its own shell, with the same algebra (argued, not computed).

\emph{(iv).} The check builds the one-particle matrices, verifies nilpotency and $KQ_1 = Q_1^{\mathsf T}K$ for the Krein Gram $K$, solves the linear system for an odd $R_1$ with $Q_1R_1 + R_1Q_1 = P_{q}$ ($P_{q}$ computed from the adapted basis), and verifies the Krein orthogonality and Gram of the physical span. On $S(V)$ it applies $d\Gamma$ monomial by monomial on every sector of fixed total number $T$ (sizes $1, 6, 19, 44, 85$ for the gluon and $1, 7, 26, 70$ for the massive vector) and confirms $Q_{B}^2 = 0$, $\langle u, Q_{B}w\rangle_K = \langle Q_{B}u, w\rangle_K$ and $Q_{B}R + RQ_{B} = d\Gamma(P_{q})$ on each basis monomial. It computes the ranks of $Q_{B}$ from ghost number $-1$ to $0$ and from $0$ to $1$, the cohomology dimensions ($\binom{T+1}{1}$, $\binom{T+2}{2}$ and $\binom{T+3}{3}$, the symmetric powers of spaces of dimension $2$, $3$ and $4$) and the Krein signature on $\ker Q_{B}$, $(\dim H, 0, \dim\operatorname{im}Q_{B})$. For the gluon the dimensions $1, 2, 3, 4$ at $T \le 3$ are those of the stored result.

\emph{(v).} The gauge algebra has twelve generators (Theorem \claimref{standard-model.gauge-algebra-minimal-massive-composite}). On a hypercharge-$\tfrac12$ doublet of \claimref{synthesis.linear-rung-field-has-a-neutral-normal-pair-in-the-scalon-sum} they act as an anti-Hermitian basis of $\mathfrak{su}(2)$, as zero for colour, and as $\tfrac{1}{2}I_{4}$ for the $\mathfrak{u}(1)$ (the unit of this parenthetical $2 \times 2$ reading is the doublet's complex structure, the field theory's unit $I_{4}$, numbers read in $\mathrm{span}\{1, I_{4}\}$). The vacuum is input: both doublets neutral and aligned, $(0, 3)$ and $(0, 4)$ in the check (any nonzero values give the same ranks). The check computes the real $8 \times 12$ Goldstone map $a \mapsto (T_a\phi_1, T_a\phi_2)$: rank $3$, kernel of dimension $9$, of which one vector lies in the electroweak block with nonzero components along both the diagonal $\mathfrak{su}(2)$ generator and hypercharge (the unbroken $\mathfrak{u}(1)$). The vector mass matrix is the Gram matrix of this map, so three vectors are massive and nine massless, and the three image directions are the Goldstone modes eaten, leaving $8 - 3 = 5$ physical scalars. It then forms the direct sum of nine gluon-type and three massive one-particle spaces ($75$ modes) and confirms, block by block, Krein self-adjointness, $Q_1^2 = 0$ and $Q_1R_1 + R_1Q_1 = P_{q}$; $P_{q}$ has rank $48$, and the Krein Gram matrix of the $27$ physical basis vectors has signature $(27, 0, 0)$. Giving the timelike polarisation the positive form instead breaks the self-adjointness and fails the check.

\emph{What is standard and what is not.} The quartet mechanism and the homotopy argument are standard (Fact \claimref{synthesis.imported-brst-kugo-ojima-strocchi-and-gribov}(ii)); new are the exact checks on the store's representations and the extension of the stored three-quanta computation to every number of quanta, for gluons and massive vectors. Clauses (K1) and (K2) of the Krein-form premise (KF, \claimref{synthesis.krein-form-premise-indefinite-inner-product-for-gauge-fields}) are hypotheses here: that the polarisation state space carries the Gupta-Bleuler form and the ghost pair the pairing $\langle c, \bar c\rangle_K = 1$; neither is derived. The breaking pattern of (v) takes the vacuum as input. The construction is perturbative: beyond perturbation theory gauge fixing has Gribov copies (Fact \claimref{synthesis.imported-brst-kugo-ojima-strocchi-and-gribov}(iv); computed lattice copies in \claimref{synthesis.non-abelian-gauge-fixing-has-gribov-copies-and-the-weak-rotor-lattice-with-doublets-has-a-perimeter-law}). The interacting statement, on the asymptotic fields of the interacting theory, needs Kugo and Ojima's hypotheses and is not proved here.
